\section{Modelo em software}

A organização lógica atual do modelo Python é:

\begin{verbatim}
model/
    vector_state.py
    vector_isa.py
    memory.py
    golden_model.py
    instruction.py
    executor.py
    trace.py
    strip_mining.py
    inspect_trace.py
    tests/
\end{verbatim}

Responsabilidades principais:

\begin{description}
    \item[\texttt{vector\_state.py}]
    Estado arquitetural, configuração SEW/LMUL e acesso lógico aos
    registradores vetoriais.

    \item[\texttt{vector\_isa.py}]
    Semântica funcional das operações vetoriais.

    \item[\texttt{memory.py}]
    Memória endereçada por byte usada pelo modelo.

    \item[\texttt{golden\_model.py}]
    API de alto nível para o estado, operações e memória.

    \item[\texttt{instruction.py}]
    Representação e validação de instruções do programa.

    \item[\texttt{executor.py}]
    Interpretação sequencial do programa e geração do trace.

    \item[\texttt{trace.py}]
    Registro e serialização dos efeitos de execução.

    \item[\texttt{strip\_mining.py}]
    Geração de programas de processamento por blocos com residual.

    \item[\texttt{inspect\_trace.py}]
    Execução de um caso demonstrativo e impressão do trace.
\end{description}